\documentclass[11pt]{article}

\usepackage[margin=1in]{geometry}
\usepackage{amsmath,amssymb,amsthm}
\usepackage{enumitem}
\usepackage{hyperref}

\title{A Minimal Mathematical Formulation of ReAct\\\large Lecture Note}
\author{Qining Shen}
\date{September 2026}

\begin{document}

\maketitle

\section{Purpose and Scope}

ReAct interleaves language-model reasoning with actions that interact with an external environment \cite{yao2023react}. This note gives that interaction a minimal mathematical representation.

The goal is descriptive rather than algorithmic. We want to make precise what information is available to the agent, how reasoning changes the agent's internal information state, how external actions interact with the environment, and how new observations enter subsequent decisions. We do \emph{not} assume a particular reinforcement-learning algorithm, a separate routing policy, or a multi-agent architecture.

\section{Basic Setup}

Let \(q\) denote the task or query. The external environment has
\[
\mathcal{X}
\quad \text{(environment states)},\qquad
\mathcal{A}
\quad \text{(external actions)},\qquad
\mathcal{O}
\quad \text{(observations)}.
\]

We allow the environment state \(x_t\in\mathcal{X}\) to be latent. An external action \(a_t\in\mathcal{A}\) changes or queries the environment according to
\begin{align}
    x_{t+1} &\sim P(\cdot\mid x_t,a_t), \label{eq:transition}\\
    o_{t+1} &\sim O(\cdot\mid x_{t+1}), \label{eq:observation}
\end{align}
where \(P\) is the transition kernel and \(O\) is the observation kernel. In deterministic tool environments, these kernels may be deterministic.

A search call, calculator call, API request, or game move can all be represented as different choices of \(a_t\). In particular, retrieval does not require a separate mathematical object at this stage; it is one possible external action.

\section{Interaction History}

At round \(t\), define the interaction history
\begin{equation}
    h_t
    =
    \bigl(
    q,
    r_0,a_0,o_1,
    r_1,a_1,o_2,
    \ldots,
    r_{t-1},a_{t-1},o_t
    \bigr),
    \label{eq:history}
\end{equation}
where
\begin{itemize}
    \item \(r_k\) is the reasoning generated by the language model at round \(k\),
    \item \(a_k\) is the external action taken at round \(k\), and
    \item \(o_{k+1}\) is the observation returned after that action.
\end{itemize}

Thus, \(h_t\) contains the information accumulated by the agent before it generates the next reasoning step. Initially,
\[
h_0=(q).
\]

We keep the full history rather than assuming a compact state representation in advance. This avoids imposing a Markov representation before a sufficient state has been identified.

\section{Internal Environment and Embedded MDP View}

To distinguish internal reasoning from interaction with the external environment, introduce an internal information state
\[
c_t\in\mathcal{C}.
\]
We write
\begin{equation}
    c_t=\psi(h_t),
    \label{eq:internal-state}
\end{equation}
where \(\psi\) represents the information from the interaction history that is available for the next computation.

This notation is not meant to introduce an additional latent \emph{mental state} of the language model. In the minimal formulation, we simply take
\begin{equation}
    c_t=h_t.
    \label{eq:internal-state-minimal}
\end{equation}
A more compact representation can be introduced later if one can identify a sufficient summary of the history.

Let \(\mathcal{R}\) denote the space of reasoning steps. A reasoning step \(r_t\in\mathcal{R}\) acts on the internal state according to
\begin{equation}
    c_t^{+}
    \sim
    P_{\mathrm{int}}(\cdot\mid c_t,r_t),
    \label{eq:internal-transition}
\end{equation}
where \(P_{\mathrm{int}}\) is the internal transition rule. In the minimal textual realization,
\begin{equation}
    c_t=h_t,
    \qquad
    c_t^{+}=h_t\oplus r_t,
    \label{eq:internal-transition-minimal}
\end{equation}
with \(\oplus\) denoting concatenation.

The important distinction is that this internal transition changes the information available to the agent without directly changing the external environment state \(x_t\). When \(c_t\) is sufficient for the internal transition dynamics, \((\mathcal{C},\mathcal{R},P_{\mathrm{int}})\) has the state--action--transition structure of an embedded MDP. We do not assign a separate intrinsic reward to this internal process; the usefulness of reasoning is evaluated through the downstream task reward of the complete ReAct trajectory.

\section{One ReAct Round}

A single ReAct round couples one internal reasoning transition with one external interaction.

\subsection{Reasoning}

Conditioned on the current history, the language model generates a reasoning step
\begin{equation}
    r_t
    \sim
    \pi^{\mathrm{R}}_{\theta}(\cdot\mid h_t).
    \label{eq:reason}
\end{equation}

The reasoning step updates the internal state according to \eqref{eq:internal-transition}, while the external state remains unchanged:
\[
(c_t,x_t)
\;\xrightarrow{\;r_t\;}\;
(c_t^{+},x_t).
\]

\subsection{External action}

Conditioned on the history and the newly generated reasoning, the model produces an external action
\begin{equation}
    a_t
    \sim
    \pi^{\mathrm{A}}_{\theta}(\cdot\mid h_t,r_t).
    \label{eq:action}
\end{equation}

The notation \(\pi^{\mathrm{R}}_{\theta}\) and \(\pi^{\mathrm{A}}_{\theta}\) does not imply two different models. They can be viewed as two conditional components of the same autoregressive language model. Equivalently,
\begin{equation}
    \pi_{\theta}(r_t,a_t\mid h_t)
    =
    \pi^{\mathrm{R}}_{\theta}(r_t\mid h_t)
    \pi^{\mathrm{A}}_{\theta}(a_t\mid h_t,r_t).
    \label{eq:factorization}
\end{equation}

After \(a_t\) is executed, the external environment evolves according to \eqref{eq:transition}--\eqref{eq:observation}. The returned observation is then incorporated into the history:
\begin{equation}
    h_{t+1}
    =
    h_t\oplus(r_t,a_t,o_{t+1}),
    \label{eq:update}
\end{equation}
and the next internal state is
\begin{equation}
    c_{t+1}
    =
    \psi(h_{t+1}).
    \label{eq:next-internal-state}
\end{equation}

Under the minimal choice \(c_t=h_t\), this reduces exactly to the original history update.

The complete round can therefore be summarized as
\[
\boxed{
(c_t,x_t)
\;\xrightarrow{\;r_t\;}\;
(c_t^{+},x_t)
\;\xrightarrow{\;a_t,\;o_{t+1}\;}\;
(c_{t+1},x_{t+1}).
}
\]

Equivalently, at the level of the observable interaction,
\[
\boxed{
h_t
\;\longrightarrow\;
r_t
\;\longrightarrow\;
a_t
\;\longrightarrow\;
o_{t+1}
\;\longrightarrow\;
h_{t+1}.
}
\]

The second display is the basic ReAct loop; the first makes explicit the different roles of internal reasoning and external interaction.

\section{Termination and Task Objective}

The interaction ends when the agent produces a finishing action,
\[
a_T=\mathrm{Finish}(\hat y),
\]
where \(\hat y\) is the final answer or solution.

Let
\[
\tau
=
(q,r_0,a_0,o_1,\ldots,r_T,a_T)
\]
denote the complete trajectory, and let \(R(\tau)\) be its task reward. For a task with a verifiable answer \(y^\star\), a simple example is
\[
R(\tau)
=
\mathbf{1}\{\hat y=y^\star\}.
\]

If one later chooses to optimize the policy parameters, the natural task-level objective is
\begin{equation}
    J(\theta)
    =
    \mathbb{E}_{\tau\sim p_{\theta}}
    \left[
        R(\tau)
    \right],
    \label{eq:objective}
\end{equation}
where \(p_\theta\) is the trajectory distribution induced jointly by the language-model policy and the environment.

Equation \eqref{eq:objective} does \emph{not} imply that ReAct must be trained with reinforcement learning. ReAct can be used with a fixed pretrained model through prompting alone. The objective only specifies what would be optimized if learning were introduced later.

\section{Sequential-Decision Interpretation}

The formulation distinguishes two coupled state spaces:
\[
x_t\in\mathcal{X}
\qquad\text{and}\qquad
c_t\in\mathcal{C}.
\]
The external state \(x_t\) describes the environment with which the agent interacts, while the internal state \(c_t\) describes the information currently available to the agent for computation.

Reasoning and acting affect these states differently:
\begin{align*}
    r_t &: \text{acts primarily on the internal state } c_t,\\
    a_t &: \text{acts on or queries the external state } x_t.
\end{align*}
The resulting observation then carries information from the external environment back into the next internal state.

This gives the following coupled process:
\[
\text{internal computation}
\;\longrightarrow\;
\text{external interaction}
\;\longrightarrow\;
\text{new observation}
\;\longrightarrow\;
\text{internal computation}.
\]

If the pair \((c_t,x_t)\) is sufficient for the corresponding transition dynamics, the full process admits a Markov representation on the joint state space. When \(x_t\) is latent, the same interaction is naturally partially observed from the agent's perspective.

For the minimal formulation in this note, however, no compact Markov state needs to be identified. Taking \(c_t=h_t\) preserves the original history-dependent description exactly. The internal MDP view is therefore best understood as an embedded component of the full ReAct interaction, not as an independent decision problem.

\section{Boundary of the Present Formulation}

This note intentionally stops at the single-agent ReAct interaction above. In particular, it does not yet introduce
\begin{itemize}
    \item a separate high-level controller deciding when to reason or act,
    \item a separate intrinsic reward for internal reasoning,
    \item an explicit cost for tokens, tool calls, or latency,
    \item a multi-agent communication or delegation mechanism, or
    \item a particular learning algorithm.
\end{itemize}

These are possible extensions, but they should be added only after the basic formulation is agreed upon.

\begin{thebibliography}{9}

\bibitem{yao2023react}
Shunyu Yao, Jeffrey Zhao, Dian Yu, Nan Du, Izhak Shafran, Karthik Narasimhan, and Yuan Cao.
\newblock ReAct: Synergizing Reasoning and Acting in Language Models.
\newblock \emph{International Conference on Learning Representations (ICLR)}, 2023.

\end{thebibliography}

\end{document}
